\author{OTS}
\documentclass[14pt]{../../inc/mybib}

\setincpath{../../inc/}

\usepackage{bible_style}
\usepackage{bibeltext}
\usepackage{markierungen}
\graphicspath{{../../assets/images/}}
\usepackage{header}

\newcommand{\Name}{Paul}

% ensure scrlayer-scrpage has sufficient footheight
\setlength{\footheight}{30.4pt}

\begin{document}
\begin{bibelbox}{SCHL}{Kol}{3:15-17}
    Lasst das Wort des Christus reichlich in euch wohnen in aller Weisheit; lehrt und ermahnt einander und singt mit Psalmen und Lobgesängen und geistlichen Liedern dem Herrn lieblich in eurem Herzen. Und was immer ihr tut in Wort und Werk, das tut alles im Namen des Herrn Jesus und dankt Gott, dem Vater, durch ihn.
\end{bibelbox}
\section{Begrüssung}
Ich möchte auch euch alle recht herzlich zum Gottesdienst begrüssen. Schön, dass ihr gekommen seid, um unserem \herr N zu danken, ihn zu loben und zu preisen.

Auch dich \Name{} will ich herzlich begrüssen. Du bist heute das erste Mal hier bei uns und wir freuen uns sehr dich hier begrüssen zu können.
\begin{itemize}
    \item[] \beten[b,p]{Beten zur Begrüssung}
    \item[] Und anschliessend singen wir zusammen das Lied
    \item[] \lied[b]{Christus ist mein ganzer Halt}
\end{itemize}

\section{Informationen}
\begin{itemize}
    \item \info[b,p]{Bibel und Gebetsabend:} Do, 04. Juni 2026 20:00 Uhr Bibel und Gebetsabend mit Markus 8:1-9 Philipp Ottenburg
    \item \info[b,p]{Nächster Gottesdienst:} So, 07. Juni 2026 Fredy Peter
    \item \info[b,p]{Kollekte:} Die Kollekte wird hinten in der grünen Box gesammelt und wird vollumfänglich für den Bau dieser Gemeinde eingesetzt.
\end{itemize}

\section{ Input }
\begin{spacing}{1.5}
    Massada selber kommt in der Bibel nicht vor. Die Festtung wurde erst von Herodes dem Grossen gebaut. Beim Totem Meer ist das anders, dieses kommt mehrfach in der Bibel vor. Es wird oft \enquote{Meer der Arava} genannt. So hiess die Grabenregion in der das Meer liegt. Ich will euch etwas über das erzählen, was in der Region zu Abrahamszeiten passiert ist. Das erste Mal kommt das dieses Meer in \bibleverse{IMos}(14:3) vor dort steht:
    \begin{bibelbox}{SCHL}{IMos}{14:3}
    Diese verbündeten sich im Tal Sidim, wo das Salzmeer ist (oder sein wird).
    \end{bibelbox}
    Aber lesen wir zuerst mal was denn dort genau passiert ist.
    \bibleverse{IMos}(14:1-16);

    Es geht also um die Entführung Lots und wie Abram ihn gerettet hat. Aber wie kam es dazu, dass Abram und Lot in dieser Gegend waren? Gott hat Abram in AI auserkoren der Stammvater seines Volkes zu werden. 
    \begin{bibelbox}{SCHL}{IMos}{12:1}
    Der \herr{} aber hatte zu Abram (noch heisst er nicht Abraham) gesprochen: Geh hinaus aus deinem Land und aus deiner Verwandschaft und aus dem Haus deines Vaters in das Land, das ich dir zeigen werde.
    \end{bibelbox}
    Abram zog los und nahm seinen Neffen Lot mit. Beide zogen zuerst nach Ägypten und kamen dann über Ägypten in den Negev die südlich vom Toten Meer liegt. 
    Mit der Zeit wurden ihre Herden aber so gross, dass es beim weiden des Viehs, viele Probleme gab und dadurch auch Streit zwischen den Hirten und mit der dort ansässigen Bevölkerung gab. Zu der Zeit wohnten dort die Kaanaiter und die Pherisiter. Wenn sich Abram an die Anweisungen Gottes gehalten hätte, wäre es nicht zu diesen Problemen gekommen. Weil Gott hat Abram in \bibleverse{IMos}(12:1) ausdrücklich gesagt, dass er alleine ausziehen soll ohne Verwandschaft.
    
    Es ist heute als Bibelleser einfach, ihm dies vorzuhalte. Abram war alleine in einem fremden Land und er hatte die Verantwortung für seinen Neffen. Aber, er hat nun mal gegen die Anweisung Gottes gehandelt.

    Zu der Zeit hat Gott Abraham direkt angesprochen. Wie weiss ich nicht, aber so, dass Abram darauf hin direkt loszog. Wir werden nicht mehr direkt von Gott angesprochen. Oft spricht er zu uns durch das Gebet und wir wissen, das kommt von Gott. Aber auch die Schrift zeigt uns den Willen Gottes. Aber wie oft, weichen wir von seinen Weisungen ab? Als ich gefragt wurde, hier in dieser Gemeinde mitzuarbeiten, hatte ich eigentlich schon recht früh die Gewissheit, dass es das ist, was Gott will. Trotzdem war ich sehr unsicher ob es wirklich von Gott ist. Darum habe ich ein paar Vertraute gefragt, was ich denn nun machen soll, um eine Bestätigung zu erhalten. Intressant ist, dass alle ausser Janina und Fredy, mir abgeraten haben. Aber zurück zu Abram.

    Abram nahm also Lot mit und stand jetzt vor dem Problem, dass die Herden zu gross wurden. Die logische Schlussvolgerung daraus war, sich zu trennen. Vielleicht hatte Abram ein schlechtes Gewissen, aber er überliess Lot die erste Wahl wohin er sich wenden will. 
    
    Lot guckte sich um und nimmt gleich das beste Stück. So lesen wir in
    \begin{bibelbox}{SCHL}{IMos}{13:10}
    Da hob Lot seine Augen auf und sah die ganze Jordanaue; denn sie war überall bewässert, wie der Garten des \herr N, wie das Land Ägypten, bis nach Zoar hinab, bevor der \herr{} Sodom und Gomaorra zerstörte.
    \end{bibelbox}
    Hier steht also, dass das Tote Meer bei der Zerstörung von Sodom und Gomorra entstand.
    Lot vergleicht diesen Teil mit dem Garten des \herr N, also mit dem Garten Eden. Lot sah darin sein neues Paradies. So wie es hier steht, lag wohl auch Sodom und Gomorra in dieser Region. Heute wird Sodom im Süden vom Toten Meer angesiedelt. Als Janina und ich von 6 Jahren in Israel waren, haben wir ein Tell, ein Hügel der aus den Überresten einer Stadt besteht, angeschaut. Dieser Hügel wurde gerade frisch ausgegraben und hat alle anzeichen, dass diese Stadt durch eine Brandkatastrophe zerstört wurde. Diese Stadt liegt nördlich vom Toten Meer. 

    Das ist jetzt also die Ausgangslage. Wenn wir weiterlesen erfahren wir, dass Lot immer näher zu den Städten rückte und wohl allmählich die Landwirtschaft aufgab. Noch war er aber mit seinen Zelten vor den Städten Sodoms. Dann steht dort noch der Satz
    \begin{bibelbox}{SCHL}{IMos}{13:13}
    Aber die Leute von Sodom waren sehr böse und sündigten schlimm gegen Gott.
    \end{bibelbox}

    Als \betonung[b,p]{wir} uns entschieden haben Jesus zu folgen, haben wir uns auch entschieden, uns dieser sündigen Welt zu entfernen. Das heisst nicht, dass wir in einem Kloster leben müssen, sondern dass wir uns von den Sünden dieser Welt fern halten sollen. Die Verlockungen der Welt sind so gross, dass es gefährlich ist, sich zu nahe an diesen Verlockungen aufzuhalten. Disco, Nachtclubs, aber auch, \betonung[b,p]{was} schaue ich im Internet, im Fernseh oder auf dem Handy an. Nicht alles was da auftaucht gefällt Gott. 
    
    Wir lesen nicht, dass Lot in Sodom mitgemacht hat, wohl nicht, weil Gott hat ihn ja, bevor er Sodom und Gomorra vernichtet hat, vor dem Untergang gerettet. Sein Aufenthalt dort, hat ihn aber viel gekostet.  
    
    Es ist darum sicherer einen gewissen Abstand von diesen sündigen Verführungen und Verführer dieser Welt zu haben, um nicht mit hineingezogen zu werden.

    Lot ist also in der Nähe dieser sündigen Stadt und wurde dadurch mit in einen Krieg verwickelt. Wir haben in unserem Text in den Versen 1 und 2 gelesen, dass sich 4 Könige gegen 5 Könige zum Krieg versammelt hatten. In der Region verlief die Handelsstrasse. Sodom und Gomorra, waren seit 14 Jahren dem König Kedor-Laomer untertan. Die Städte wollten frei werden und da Kedor-Laomer auf einem Kriegszug war, dachten sie, er sei geschwächt und griffen diesen König an. Wahrscheinlich, hat Lot sich in seinen Zelten vor Sodom aufgehalten und ist so zwischen die Fronten geraten und wurde überfallen und entführt.
    
    Das passiert auch uns, wenn wir zu nahe an der Sünde sind. Wir können zwischen die Fronten geraten. Noch schlimmer wenn wir uns mit ihr verbinden. Es ist dann nur noch ein kleiner Schritt zurück in die Sünde.  Wie heisst es so schön? \enquote{Mitgegangen, Mitgehangen}. Nach der Bekehrung hat man vielleicht aufgehört mit der Internetpronographie, aber man hat noch ein paar Seiten auf den PC nicht gelöscht oder noch Bilder auf dem Computer und irgendwann ist man wieder in diesem Sumpf gefangen. Oder man ist wieder mit den ehemalige Saufkumpanen unterwegs. Auch wenn man gute Absichten hat und am Anfang nur Wasser trinkt, es ist gefährlich sich so nahe an der Sünde zu bewegen.

    Lot wurde entführt und dies wurde Abram berichtet. Abram nahm 318 seiner erprobten Knechte und jagte den Entführer nach und befreite Lot. Die Gegner waren schon auf dem Heimweg als Abram die Meldung bekam. Es sagt sich so schnell, aber Abram war hier unten bei Hebron und ist dann losgezogen  bis nach Dan hoch, hier ganz oben und erst hier, konnte er Lot befreien. Abram hat Lot nicht im Stich gelassen. Er konnte seinen Neffen befreien, ihn und sein ganzer Hausrat.

    Wenn auch wir wieder zurück in eine Sünde fallen, dürfen auch wir immer wieder unseren Herrn anrufen. Ich rede hier nicht von unseren täglichen Sünden, die sind nicht harmloser und sollten immer direkt vor den Herrn gebracht werden und busse getan werden.

    Aber wenn wir uns von ihm getrennt haben und in einer Sünde verharren, er wird uns nicht im Sumpf stecken lassen. Er hilft uns. So auch in seinem Gleichnissen vom Verlorenen Sohn, oder dem Schaf das verloren ging. Er hilft uns immer wieder aus diesem Sumpf der Sünde rauszukommen. Er zeigt uns einen Weg, vieleicht durch die Geschwister aus der Gemeinde oder auch durch eine professionelle Hilfe. Wir dürfen das nicht auf die leichte Schulter nehmen. Wir müssen ihn darum bitten und auch danach handeln. Lot wurde dieses mal von Abram gerettet, ist aber direkt wieder zurück nach Sodom und vermutlich wohnte er in der Stadt, weil er seine Zelte und Felder im Krieg verloren hat.

    Nach der Rückkehr, nachdem wir die Sünde wieder in den Griff bekommen haben, sollte wir als erstes eine Distanz zu dieser Sünde schaffen. Wegen meiner Trinkerei muss ich mich nicht nur vom Alkohol entfernen sondern auch von dem damaligen Umfeld. Nach jedem Rückfall wird es immer schwieriger da wieder rauszukommen. Jedes mal bleiben Narben zurück. Irgendwann sagt man sich: \enquote{Ich schaffe das sowieso nicht!} lässt sich gehen und fällt in Depressionen.  Wir werden das Heil nicht verlieren. Lot wurde gerettet, aber er hat alles verloren, er hatte nur noch seine zwei Töchter sonst war alles weg. Wie schon gesagt, das Heil können wir nicht verlieren, aber hier im Leben die Segnungen unseres Herrn die er uns versprochen hat, und wir können den hl. Geist betrüben und dämpfen. Darum direkt mit jeder Sünde vor den Herrn und busse tun.

    \begin{block}[Tote Meer]
    \end{block}
    \begin{block}[Massada]
    \end{block}
    \lied[b]{Alle Herren dieser Welt}.
\end{spacing}
\section{Predigt}
\textbf{Nach der Predigt}\newline
Danken für die Predigt.
% \section{Abendmahl}
% \begin{itemize}
%     \item [] Beten für das Brot
%     \item [] Beten für den Wein
% \end{itemize}
\section{Abschluss}
\begin{itemize}
    \item [] \lied[b]{Wir haben eine Hoffnung}.
    \item [] \beten[b]{Zum Abschluss}
\end{itemize}
\begin{bibelbox}{SCHL}{IIKor}{13:13}
    Die Gnade des Herrn Jesus Christus und die Liebe Gottes und die Gemeinschaft des Heiligen Geistes sei mit euch allen! Amen
\end{bibelbox}
\end{document}